\chapter{مفهوم حد و پیوستگی}

\section{تعریف شهودی حد}
در این بخش با مفهوم شهودی حد توابع آشنا می‌شویم.
اگر تابع $f(x)$ در همسایگی نقطه $x_0$ تعریف شده باشد:
$$\lim_{x \to x_0} f(x) = L$$

\begin{question}
حاصل حد $\lim_{x \to 2} (3x + 1)$ کدام است؟
\begin{choices}
\choice $5$
\choice $6$
\choice $7$ *
\choice $8$
\end{choices}
\begin{explanation}
کافی است مقدار $x = 2$ را در ضابطه تابع قرار دهیم:
$$3(2) + 1 = 6 + 1 = 7$$
\end{explanation}
\end{question}

\begin{flashcard}{قضیه فشردگی (ساندویچ)}
اگر در یک همسایگی محذوف از $x_0$ داشته باشیم $g(x) \le f(x) \le h(x)$ و $\lim g(x) = \lim h(x) = L$، آنگاه:
$$\lim_{x \to x_0} f(x) = L$$
\end{flashcard}
